\documentclass[10pt,a4paper]{amsart}
\usepackage[margin=19mm]{geometry}
\usepackage{amsmath,amssymb,mathtools}
\usepackage[scr=rsfs,cal=euler]{mathalfa}
\usepackage{xfrac}
\usepackage[unicode,hidelinks,bookmarks=false]{hyperref}
\newcommand{\ie}{i.e.\ }
\newcommand{\cf}{cf.\ }
\newcommand{\resp}{respectively\ }
\newcommand{\Subsection}{\subsection}
\providecommand{\nobreakdash}{\nobreak\hbox{-}}
\setlength{\emergencystretch}{2em}
\multlinegap0pt
\usepackage{amssymb}
\usepackage{enumerate}
\usepackage{xcolor}
\usepackage{tikz}
\newtheorem{lemma}{Lemma}
\newtheorem{theorem}{Theorem}
\title{Winning Criteria for Open Games: A Game-Theoretic Approach to Prefix Codes}
\author{Dean Kraizberg}
\date{Source: arXiv:2601.17521v2 (CC BY 4.0)}
\begin{document}
\maketitle

\noindent\emph{Source and licence.} Winning Criteria for Open Games: A Game-Theoretic Approach to Prefix Codes. Version arXiv:2601.17521v2 (CC BY 4.0), distributed by arXiv under the Creative Commons Attribution 4.0 International licence, http://creativecommons.org/licenses/by/4.0/, as recorded in the arXiv licence metadata for this exact version and shown on the abstract page https://arxiv.org/abs/2601.17521v2. Full licence notice: \texttt{licenses/arxiv-2601.17521-CC-BY-4.0.txt}.\par\smallskip

\noindent\textit{Editorial scope.} Counted unit: the article's theorem housdorf (source0.tex lines 566--570). The article's complete original proof of it is delivered with it (source0.tex lines 572--746, one \texttt{proof} environment of 175 lines). No mathematics is written, repaired or re-derived: the statement and the proof are the article's own lines, in the article's own notation, and no retained line is substituted or re-flowed.  Retained uncounted as the source-local context the proof needs: lines 557--563.  Nothing was omitted from any retained span, so the package carries no omission record.  No retained line contains a citation, so the wrapper carries no bibliography and no bibliography entry.  No retained line contains a figure environment, an image-inclusion command or drawing code, so no image is delivered and the package declares no asset dependency.  Editorial formatting only, in addition: an AMS \texttt{amsart} preamble that restates the article's own declarations and macros used by the retained lines, namely the environments \texttt{lemma}, \texttt{theorem} on separate counters, together with the standard packages those lines need.  The retained environments print in this document as Lemma 1, Theorem 1 in order of appearance, whereas the article prints the same items with the numbering of its own full document.  The complete native source of the article as distributed in the arXiv e-print for this exact version is bundled unchanged as \texttt{sources/193321/source0.tex}.
\begin{lemma} \label{lemma: Housdorf and sum 1/2}
Let $T=\mathcal A^{\leq \mathbb N}$ and let $W=\bigcup_{p\in Z}[T_p]$ be an open set, with $Z \subseteq T$ being a finite collection of even positions.  
Then $\dim_H(Z_{\mathrm{concat}})<\tfrac12$, if and only if
\[
\sum_{p\in Z} \frac{1}{|\mathcal A|^{ \mathrm{len}(p)/2 }} < 1.
\]
\end{lemma}

\begin{theorem}\label{theorem Fishman game housdorf}
Let \(T=\mathcal{A}^{\le \mathbb{N}}\) and let \(S\subseteq [T]\) be a Borel set. 
If there exist integers \(k\ge 1\) and a set \(A\subseteq \mathcal{A}^k\) with \(|A|\ge |\mathcal{A}|^{k d}\) and \(A_{\mathrm{concat}} \subseteq S\),
then \(\dim_H(S)\ge d\).
\end{theorem}

\begin{proof}[Proof of Lemma~\ref{lemma: Housdorf and sum 1/2}]
We recall the definition of Hausdorff dimension.

\smallskip

Let \(\delta>0\). A \(\delta\)-cover of \(Z_{\mathrm{concat}}\) is a countable collection \(\{U_i\}_{i\in I}\) such that \(Z_{\mathrm{concat}}\subseteq\bigcup_{i\in I}U_i\) and \(\operatorname{diam}(U_i)<\delta\) for every \(i\in I\). The \(d\)-dimensional Hausdorff measure of \(Z_{\mathrm{concat}}\) is
\[
H^d(Z_{\mathrm{concat}})
=
\liminf_{\delta\to 0}
\left\{
\sum_{i\in I}(\operatorname{diam} U_i)^d
\;\middle|\;
\{U_i\}_{i\in I}\text{ is a }\delta\text{-cover of }Z_{\mathrm{concat}}
\right\},
\]
and the Hausdorff dimension is
\[
\dim_H(Z_{\mathrm{concat}})
=
\inf\{d\ge 0:\; H^d(Z_{\mathrm{concat}})=0\}.
\]

Let \(2n=\max_{p\in Z}\operatorname{len}(p)\), and fix \(p\in Z\). We replace \(p\) by a uniformly expanded family of extensions of length \(2n\), constructed by inserting a fixed symbol \(a_0\in\mathcal A\) in every even position and allowing every symbol of \(\mathcal A\) to appear in the intervening odd positions.

Recall that we have assumed every \(p\in Z\) has even length. Define
\[
\operatorname{Extend}_n(p)
:=
\Bigl\{
p \circ \langle a_0,a_1\rangle \circ \langle a_0,a_2\rangle \circ \cdots \circ
\langle a_0,a_{\,n-(\operatorname{len}(p)/2)}\rangle
:
a_1,\dots,a_{\,n-(\operatorname{len}(p)/2)}\in\mathcal A
\Bigr\}.
\]
Every \(q\in\operatorname{Extend}_n(p)\) satisfies \(\operatorname{len}(q)=2n\), and
\[
\bigl|\operatorname{Extend}_n(p)\bigr|
=|\mathcal A|^{\,n-(\operatorname{len}(p)/2)}.
\]

Set
\[
W_n := \left( \bigcup_{p\in Z} \operatorname{Extend}_n(p)\right)_{\mathrm{concat}}.
\]
We claim that if \(\dim_H(Z_{\mathrm{concat}})<1/2\) then \(\dim_H(W_n)<1/2\). To see this, let \(\{U_i\}_{i\in I}\) be a \(\delta\)-cover of \(Z_{\mathrm{concat}}\). We may assume each \(U_i\) is a basic open set: if \(\operatorname{diam}(U_i)=0\) then \(U_i\) contains a single point and can be replaced by a basic open set of arbitrarily small diameter; moreover,
\[
\operatorname{diam}(U_i)=|\mathcal A|^{-\min\{k : \exists x,y\in U_i,\; x_k\neq y_k\}},
\]
so replacing \(U_i\) by a containing basic open set does not increase its diameter.

Since the infimum defining \(H^d\) may be taken over basic open covers, we may further assume each \(U_i\) has the form \([T_p]\), where \(p\) is a prefix of a concatenation \(p_1\circ\cdots\circ p_k\) with \(p_1,\dots,p_k\in Z\). For such a basic open set \(U_i=[T_p]\) with
\(p=p_1\circ\cdots\circ p_{k-1}\circ p'\) and \(p'\preceq p_k\in Z\), define
\[
\operatorname{Extend}_n(U_i)
:=\bigl\{[T_q]:
q=q_1\circ\cdots\circ q_{k-1}\circ p',\;
q_j\in\operatorname{Extend}_n(p_j)\ \text{for }1\le j\le k-1\bigr\}.
\]
Then \(\bigcup_{i\in I}\operatorname{Extend}_n(U_i)\) is an open cover of \(W_n\).

Because \(\dim_H(Z_{\mathrm{concat}})<1/2\), there exists \(\epsilon_0>0\) such that \(H^{1/2-\epsilon_0}(Z_{\mathrm{concat}})=0\). Put \(\epsilon=\epsilon_0/(2n)\). We claim that
\[
\sum_{i\in I}\sum_{U'\in\operatorname{Extend}_n(U_i)}(\operatorname{diam} U')^{\,1/2-\epsilon}
\le
\sum_{i\in I}(\operatorname{diam} U_i)^{\,1/2-\epsilon_0}.
\]
To verify this inequality, fix \(U_i=[T_p]\) with
\(p=p_1\circ\cdots\circ p_{k-1}\circ p'\) as above. Then
\[
\sum_{U'\in\operatorname{Extend}_n(U_i)}(\operatorname{diam} U')^{\,1/2-\epsilon}
=
\sum_{\substack{q_j\in\operatorname{Extend}_n(p_j)\\ j=1,\dots,k-1}}
\biggl(\frac{1}{|\mathcal A|}\biggr)^{(2n(k-1)+\operatorname{len}(p'))(1/2-\epsilon)}.
\]
Since there are \(|\mathcal A|^{\,n-\operatorname{len}(p_j)/2}\) choices for each \(q_j\), the preceding display equals
\[
\biggl(\frac{1}{|\mathcal A|}\biggr)^{(2n(k-1)+\operatorname{len}(p'))(1/2-\epsilon)}
\prod_{j=1}^{k-1} |\mathcal A|^{\,n-\operatorname{len}(p_j)/2}.
\]
A short rearrangement yields
\[
\sum_{U'\in\operatorname{Extend}_n(U_i)}(\operatorname{diam} U')^{\,1/2-\epsilon}
=
|\mathcal A|^{\,\epsilon(2n(k-1)+\operatorname{len}(p'))}\cdot
\bigl(\operatorname{diam} U_i\bigr)^{1/2}.
\]
Since \(\epsilon(2n(k-1)+\operatorname{len}(p'))\le \epsilon_0\operatorname{len}(p)\) and \((\operatorname{diam} U_i)^{\epsilon_0} =  |\mathcal{A}|^{-\epsilon_0 \operatorname{len}(p)} \), we obtain
\[
\sum_{U'\in\operatorname{Extend}_n(U_i)}(\operatorname{diam} U')^{\,1/2-\epsilon}
\le (\operatorname{diam} U_i)^{1/2-\epsilon_0},
\]
as claimed.

It follows that
\[
H^{1/2-\epsilon}(W_n)
\le
\liminf_{\delta\to 0}
\Biggl\{
\sum_{i\in I}\sum_{U'\in\operatorname{Extend}_n(U_i)}(\operatorname{diam} U')^{\,1/2-\epsilon}
\;\Bigm|\;
\{U_i\}_{i\in I}\ \text{is an open }\delta\text{-cover of }Z_{\mathrm{concat}}
\Biggr\}
\]
\[
\le
\liminf_{\delta\to 0}
\Biggl\{
\sum_{i\in I}(\operatorname{diam} U_i)^{\,1/2-\epsilon_0}
\;\Bigm|\;
\{U_i\}_{i\in I}\ \text{is an open }\delta\text{-cover of }Z_{\mathrm{concat}}
\Biggr\}
=0.
\]
Hence \(\dim_H(W_n)<1/2\).

\medskip

Denote
\[
A:=\bigcup_{p\in Z}\operatorname{Extend}_n(p)\subseteq\mathcal A^{2n}.
\]
Note that
\[
|A|=\sum_{p\in Z} |\mathcal A|^{\,n-(\operatorname{len}(p)/2)}.
\]
Since \(A_{\mathrm{concat}}\subseteq W_n \) and \(\dim_H(W_n)<1/2\), Theorem~\ref{theorem Fishman game housdorf} implies \(|A|<|\mathcal A|^n\). Therefore
\[
\sum_{p\in Z} \frac{1}{|\mathcal A|^{\,\operatorname{len}(p)/2}}<1,
\]
which is the desired inequality.

\medskip

For the converse, assume
\[
\sum_{p\in Z}\biggl(\frac{1}{|\mathcal A|}\biggr)^{d\cdot\operatorname{len}(p)}<1
\quad\text{for }d=\tfrac12.
\]
By continuity in the exponent \(d\), there exists \(\epsilon>0\) such that
\[
\sum_{p\in Z}\biggl(\frac{1}{|\mathcal A|}\biggr)^{(1/2-\epsilon)\operatorname{len}(p)}<1.
\]
For each \(n\ge 1\) the family
\[
\bigcup_{(p_1,\dots,p_n)\in Z^n} [T_{p_1\circ\cdots\circ p_n}]
\]
is an open cover of \(Z_{\mathrm{concat}}\), and
\[
\max_{(p_1,\dots,p_n)\in Z^n}\operatorname{diam}([T_{p_1\circ\cdots\circ p_n}])
=|\mathcal A|^{-n\cdot \min_{p\in Z}\operatorname{len}(p)} \underset{n \to \infty}{\to}0.
\]
Hence
\[
H^{1/2-\epsilon}(Z_{\mathrm{concat}})
\le
\lim_{n\to\infty}
\sum_{(p_1,\dots,p_n)\in Z^n}
\operatorname{diam}([T_{p_1\circ\cdots\circ p_n}])^{\,1/2-\epsilon}
\]
\[
=
\lim_{n\to\infty}
\sum_{(p_1,\dots,p_n)\in Z^n}
\biggl(\frac{1}{|\mathcal A|}\biggr)^{(1/2-\epsilon)\sum_{j=1}^n\operatorname{len}(p_j)}
=
\lim_{n\to\infty}
\Biggl(
\sum_{p\in Z}\biggl(\frac{1}{|\mathcal A|}\biggr)^{(1/2-\epsilon)\operatorname{len}(p)}
\Biggr)^{\!n}.
\]
By the assumption the last expression vanishes, therefore \(H^{1/2-\epsilon}(Z_{\mathrm{concat}})=0\) and \(\dim_H(Z_{\mathrm{concat}})<1/2\).
\end{proof}

\end{document}
